\section{Artifact Availability and Limits}

The released repository contains the \tool Python package, adapter template,
metric runners, sample inputs, tests, documentation, an MIT license, and a
clean-checkout verifier. The released resource is available at
\url{https://github.com/jdding/sidinspector} under tag
\texttt{sidinspector-cikm2026-resource-v0.6}. It is intended as a reusable
diagnostic resource: a user can clone the package, run the toy
diagnostic, validate a new item-to-\sid{} mapping, and obtain D1--D5 reports
without experiment caches. The released repository focuses on the reusable
interface and diagnostic implementation; selected experimental evidence is
reported in the paper.

The resource supports three modes. A downstream recommender
researcher can use the diagnostics as a pre-training triage step: artifacts
with missing items, inconsistent depth, severe aliasing, or weak
collaborative-prefix recovery are flagged for inspection before generator
training consumes GPU time. A method author can add an adapter that emits the
normalized SID-assignment table, then receive D1--D5 reports without
changing their training loop. A reader can run the clean-check verifier
against the reusable package. These signals complement
Recall@K and NDCG by making tokenizer artifacts inspectable before a full
benchmark or generator-training run.

The evidence supports a resource-interface claim over four tokenizer lines:
GRID-style RQ-KMeans, ReSID/GAOQ, LETTER, and LC-Rec. RQ-min and the controlled
stressors are reference rows for interpreting the diagnostics; CARD enters the
catalog once generated codes can be joined and validated. D2 reports mapping
aliasing; D3 uses prefix-neighborhood and reranker checks; D5 measures
structural cost; D6 is available for paired refresh mappings; D7 activates when generator
traces are supplied.

New adapters emit the same normalized tables and record feature or checkpoint
provenance. A tokenizer line enters the catalog when item-level codes pass the
validator and join to metadata or interactions; mechanism probes remain
diagnostic calibration rows. Future versions should add causal collision-harm
validation, broader adapter coverage, generator-output diagnostics, and
cross-domain or unified search/recommendation checks~\cite{li2024survey,hu2025gencdr,penha2025jointsid}.
Industrial SID deployment and
ranking studies further motivate treating these artifact properties as
engineering objects~\cite{singh2023bettersemanticids,li2026sidcoord,ju2026snapchatsid,xu2026largescalegenrec}.
